\subsection{\texorpdfstring{Topological vector spaces over \(\mathbf{C}\)}{Topological vector spaces over the complex field}}

Let E be a topological vector space over C the field of complex numbers; the topology of E is also compatible with the structure of the vector space over R, obtained by restricting the field of scalars to R. We denote by \(E_{0}\) the topological vector space on R which underlies E (I, p.~2). Note that, in \(E_{0}\) , the mapping \(x \mapsto ix\) (which is not a homothety) is an automorphism u of the topological vector space structure of \(E_{0}\) such that \(u^{2}(x) = -x\) .

Conversely, let \(\mathbf{F}\) be a topological vector space over \(\mathbf{R}\) , and suppose that there exists an automorphism \(u\) of \(\mathbf{F}\) such that \(u^2 = -1_{\mathbf{F}}(1_{\mathbf{F}}\) is the identity automorphism of \(\mathbf{F}\) . We know (A, IX, § 3, No.~2) that it is then possible to define a vector space structure on \(\mathbf{F}\) relative to \(\mathbf{C}\) writing \(\lambda x = \alpha x + \beta u(x)\) for all \(\lambda = \alpha + i\beta \in \mathbf{C}\) and all \(x \in \mathbf{F}\) . Further since the mapping \((\alpha, \beta, x) \mapsto \alpha x + \beta u(x)\) of \(\mathbf{R}^2 \times \mathbf{F}\) in \(\mathbf{F}\) is continuous the topology of \(\mathbf{F}\) is compatible with the vector space structure relative to \(\mathbf{C}\) defined above; if \(\mathbf{E}\) denotes the topological vector space on \(\mathbf{C}\) defined in this manner, then \(\mathbf{F}\) is the topological vector space on \(\mathbf{R}\) which underlies \(\mathbf{E}\) .

Remark. --- Given a topological vector space F over R, it is not always the case that there exists an automorphism u of F whose square is \(-1_{F}\) ; for example, it is not possible to define vector space structure relative to C on a vector space over R of finite odd dimension.

Let E be a topological vector space on C, and \(E_{0}\) the topological vector space on R which underlies E. Every linear variety M in E is also a linear variety in \(E_{0}\) , but the converse is false. To avoid confusion we say that a linear variety for a vector space structure relative to C (resp. relative to R) is a complex (resp. real) linear variety. A complex linear variety of finite dimension n (resp. of finite codimension n) is a real linear variety of dimension 2n (resp. of codimension 2n). In order that a real vector subspace M of E should also be a complex vector subspace, it is necessary and sufficient that \(iM \subset M\) .

Recall that, if E and F are two topological vector spaces on C, then a mapping of E in F is called C-linear (resp. R-linear) if it is a linear mapping for the vector space structures of E and of F relative to C (resp. R); every C-linear mapping is evidently R-linear but the converse is false. We say that a C-linear form on E is a complex linear form and that an R-linear form on E (i.e.~a linear form on \(E_{0}\) ) is a real linear form. If f is a complex linear form on E, it is clear that the real part g = Rf and the imaginary part h = If of f are real linear forms; further, the relation \(f(ix) = if(x)\) implies the identity \(h(x) = -g(ix)\) ; in other words we have

\[
f (x) = (\mathscr {R} f) (x) - i (\mathscr {R} f) (i x).\tag{1}
\]

Conversely, if g is a real linear form on E, then \(f(x) = g(x) - ig(ix)\) is the unique complex linear form on E such that Rf = g; and f is continuous if, and only if, g is continuous.

Now let H be a complex hyperplane in E, with the equation \(f(x) = \alpha + i\beta\) , where f is a complex linear form on E; putting \(g = Rf\) , we see that H is the intersection of two real hyperplanes \(H_{1}\) , \(H_{2}\) with equations respectively \(g(x) = \alpha\) and \(g(ix) = -\beta\) ; if H is closed, so also are \(H_{1}\) and \(H_{2}\) (I, p.~13, th. 1). Conversely let \(H_{0}\) be a homogeneous real hyperplane, with equation \(g(x) = 0\) (where g is a real linear form on E); then H, the intersection of \(H_{0}\) and \(iH_{0}\) , is a homogeneous complex hyperplane, and if f is the complex linear form such that Rf = g, then \(f(x) = 0\) is the equation of H; if \(H_{0}\) is closed then H also is closed.

Let G be a topological vector space over R and let \(\mathbf{G}_{(\mathbf{C})}\) be the vector space on C obtained from G by extending the field of scalars to C (A, II, § 5.1). Identify G as a subset of \(\mathbf{G}_{(\mathbf{C})}\) by the mapping \(x \mapsto 1 \otimes x\) . The R-linear mapping \((x, y) \mapsto x + i \cdot y\) is then a bijection of \(G \times G\) on \(\mathbf{G}_{(\mathbf{C})}\) , by means of which we transfer the product topology of \(G \times G\) to \(\mathbf{G}_{(\mathbf{C})}\) . Then \(\mathbf{G}_{(\mathbf{C})}\) with this topology is a topological vector space on C. We say that \(\mathbf{G}_{(\mathbf{C})}\) is the complexified topological vector space of G.

